% !TEX encoding = UTF-8
% !TEX Root = 0_main.tex

% \section{Simulated Verification}
% \vspace{-0.3cm}
\subsection{Simulated Verification}
% \vspace{-0.2cm}

% In this section, we examine the two approximations used in the previous section to analyze the variance.
% Specifically, as in Equation~\ref{eqn:variance_appro_appro}, we approximate $\Var[\sqrt{\frac{t}{\sum_{i=1}^t g_i^2}}]$ to the first order and assume $\psi^2(.) = \frac{1-\beta_2^t}{(1 - \beta_2)\sum_{i = 1}^t \beta_2^{t-i}g_i^2}$ subjects to the scaled inverse chi-square distribution. 
% In our analysis 
In Sections~\ref{sec:ana} and \ref{sec:fix}, we approximate $\Var[\sqrt{t/\sum_{i=1}^t g_i^2}]$ to the first order, and assume $\psi^2(.) = \frac{1-\beta_2^t}{(1 - \beta_2)\sum_{i = 1}^t \beta_2^{t-i}g_i^2}$ subjects to a scaled inverse chi-square distribution (this assumption covers the approximation from EMA to SMA).
Here, we examine these two approximations using simulations.

\noindent\textbf{First Order Approximation of $\Var[\sqrt{t/\sum_{i=1}^t g_i^2}]$.}
\label{subsec:approx}
% Here, we first derive the analytic form of $\Var[\sqrt{\frac{t}{\sum_{i=1}^t g_i^2}}]$ then compare its value to the first order approximation in Equation~\ref{eqn:variance_appro}.
% For ease of notation, we refer $\frac{t}{\sum_{i=1}^t g_i^2}$ as $x$, and mark $\frac{1}{\sigma^2}$ as $\tau^2$.
% Therefore, $x \sim \sics(t, \tau^2)$ and $p(x) = \frac{(\tau^2t/2)^{t/2}}{\Gamma(t/2)}\frac{\exp[\frac{-v\tau^2}{2x}]}{x^{1 + t/2}}$.
% We have: 
% \begin{equation}
%     \E[\sqrt{x}] = \int_{0}^{\infty} \sqrt{x} p(x) dx = \frac{\tau \sqrt{t} \,\Gamma(t/2 - 1)}{\sqrt{2} \,\Gamma(t/2)}.
%     \label{eqn:expect-sqrt-x}
% \end{equation}
% Based on Equation~\ref{eqn:inv-chi-sq} and \ref{eqn:expect-sqrt-x}, we have:
% \begin{equation}
% \Var[\sqrt{x}] = \E[x] - \E[\sqrt{x}]^2 = \tau^2 (\frac{t}{t-2} - \frac{t \,2^{2t - 5}}{\pi}\gB(\frac{t-1}{2}, \frac{t-1}{2})^2).
% \label{eqn:analytic-sqrt-var}
% \end{equation}
To compare Equations~\ref{eqn:variance_appro} and \ref{eqn:analytic-sqrt-var}, we assume $\tau = 1$ and plot their values and difference for $\nu = \{5, \cdots, 500\}$ in Figure~\ref{fig:var_approx}.
The curve of the analytic form and the first-order approximation highly overlap, and their difference is much smaller than their value.
This result verifies that our first-order approximation is very accurate.% of  $\Var[\sqrt{\frac{t}{\sum_{i=1}^t g_i^2}}]$.
% It is worth mentioning that, although we get the analytic form of $ \Var[\sqrt{\frac{t}{\sum_{i=1}^t g_i^2}}]$, it's preferable to use its first order approximation,
% % as Equation~\ref{eqn:variance_appro}, since the Equation~\ref{eqn:analytic-sqrt-var} is not numerical stable.
% which is more stable and efficient to calculate. 
% % (Mathematica\footnote{Version Number ?} fails to calculate the analytic value for $t = 1000$). 

\noindent\textbf{Scaled Inverse Chi-Square Distribution Assumption.}
\label{subsec:sma-ema}
In this paper, we assume $g_i$ accords to a Normal distribution with a zero mean. 
%Also, based on the similarity between the exponential moving average and simple moving average, 
We also assume $\psi^2(.)$ accords to the scaled inverse chi-square distribution to derive the variance of $\Var[\psi(.)]$,
based on the similarity between the exponential moving average and simple moving average.
Here, we empirically verify this assumption.
% via simulation. 

Specifically, since $g_i$ in the optimization problem may not be zero-mean, we assume its expectation is $\mu$ and sample $g_i$ from $\cN(\mu, 1)$.
Then, based on these samples, we calculate the variance of the original adaptive learning rate and the proposed rectified adaptive learning rate, \ie,  $\Var[\frac{1}{\widehat{v_t}}]$ and $\Var[\frac{r_t}{\widehat{v_t}}]$ respectively.
We set $\beta_2$ to $0.999$, the number of sampled trajectories to $5000$, the number of iterations to $6000$, and summarize the simulation results in Figure~\ref{fig:vt_var_simulation}. 
Across all six settings with different $\mu$, the adaptive learning rate has a larger variance in the first stage and the rectified adaptive learning rate has relative consistent variance. 
This verifies the reliability of our assumption. 
% At the same time, the magnitude of $\Var[\frac{1}{v_t}]$ subjects to not only the sample size, but also the mean (\ie, $\mu$).



% We show the simulation result of the variance of the adaptive learning rate by assuming $g_i \sim \cN(\mu,\sigma^2)$. In our previous analysis, we simply assume $\mu=0$ and approximate EMA with SMA. In this section, we verify that the exploding variance of the adaptive learning rate ($v_t = \sqrt{\frac{1-\beta_2^t}{(1 - \beta_2)\sum_{i = 1}^t \beta_2^{t-i}g_i^2}}$) exists for a large range of $\mu$, and our proposed rectification term can fix this issue. In the following experiments, we set $\beta_2$ as $0.999$, the number of sampled trajectories as $5000$, and the number of iterations as $6000$.





% \begin{figure}[h]
% \centering
%     \includegraphics[width=0.5\textwidth]{fig/variance_approx.pdf}
% \caption{The simulation of the analytic form, first order approximation of $\Var[\sqrt{\frac{t}{\sum_{i=1}^t g_i^2}}]$, and their difference (calculated as the absolute difference value). The y-axis is log scaled.}
%     \label{fig:var_approx}
% \end{figure}

% \begin{figure}[h]
% \centering
% \begin{tabular}{ ccc } 
%  \includegraphics[width=0.3\textwidth]{fig/simu_var/0_1_var.pdf} & 
%  \includegraphics[width=0.3\textwidth]{fig/simu_var/0001_1_var.pdf} & 
%  \includegraphics[width=0.3\textwidth]{fig/simu_var/001_1_var.pdf}  \\
%  $\mu=0$ & $\mu=0.001$  & $\mu=0.01$ \\ 
%  \includegraphics[width=0.3\textwidth]{fig/simu_var/01_1_var.pdf} & 
%  \includegraphics[width=0.3\textwidth]{fig/simu_var/1_1_var.pdf}& 
%  \includegraphics[width=0.3\textwidth]{fig/simu_var/10_1_var.pdf} \\
%  $\mu=0.1$ & $\mu=1$ & $\mu=10$ \\
% \end{tabular}

% \caption{The simulation of $\Var[\frac{1}{v_t}]$ and $\Var[\frac{c_t}{v_t}]$. The y-axis is log scaled.}
%     \label{fig:vt_var_simulation}
% \end{figure}
